\subsection*{Acknowledgements}

Penso di aver deluso qualcuno quando, dopo la mia discussione di tesi triennale, amici e parenti si sono ritrovati a fissare l’ultima pagina della mia tesi: una pagina completamente bianca. Niente ringraziamenti (o per gli amici anglofoni, \textit{acknowledgements}), e nessuna frase a effetto. 
La risposta che davo era sempre la solita: \textit{``Non ho avuto tempo''}. La verità è che non mi importava granché: la laurea triennale per me non è mai stata un traguardo o una pietra miliare per cui sentirsi appagato. La vedevo solo come una tappa intermedia, un passaggio necessario per poter andare a studiare di più e altrove, in un posto che allora ritenevo più stimolante di Bologna.

Oggi, al termine di questo percorso tra Bologna, Barcellona e Zurigo, sento che è arrivato il momento di ringraziare tutte le persone che hanno messo un tassello fondamentale lungo la strada.


{
  \setlength{\parskip}{1em}
  \setlength{\parindent}{0pt}

\vspace{3em}

A mio fratello \textbf{Samuele}, per avermi sempre protetto in silenzio e per avermi fatto capire, anche senza troppe parole, che continuo a essere il tuo fratellino. Per me rimarrai sempre Samuelona, l’uomo più intelligente del mondo.

A mia sorella \textbf{Agnese}, per avermi insegnato che la vita va presa di petto e che non bisogna avere paura di quello che si desidera essere (e per avermi insegnato a non mettere la virgola tra verbo e soggetto).

A \textbf{Nonno Francesco}, per i ciambelloni trafficati nella valigia che mi fanno mancare un po' meno casa, e per essere sempre e costantemente il mio fan numero uno.

A \textbf{Nonna Rosa}, per i 50 euro che ogni volta mi infili in tasca per comprarmi la pizza in Svizzera (anche se qui la moneta corrente è il franco). Grazie per insegnarmi a non avere paura di dire la propria, e per ricordarmi il valore della famiglia.

Alla mia \textbf{mamma}, per essere la prima a venirmi ad accogliere e l'ultima a salutarmi ogni volta che torno, non prima di avermi stipato di cibo la valigia. Per la pazienza e la gentilezza gratuita con cui ti rapporti alle persone, per insegnarmi senza imporre nulla, per la dichiarazione dei redditi, per la stima che riservi alla mia opinione, per non perdere fiducia dopo le note a scuola, per i libri del WWF letti prima di andare a letto, per non avermi mai rinfacciato niente, per la tua cucina, per le vecchie foto custodite in giro per la casa, per non chiedere mai nulla in cambio.

E al mio \textbf{babbo}, per ricordarmi il valore della famiglia, ed avermi dato ogni volta nuove possibilità (anche quando ho grippato la moto). Per esserti fatto in quattro senza volere niente in cambio, per l'interesse per me, per i tuoi bicipiti da palestrato, per i film serali, per la tua sempre professionale eleganza, per avermi insegnato passione e dedizione al lavoro, per le tue camicie anni 90, per una vita senza mai impormi niente.

A \textbf{Pelle}, per essere sempre pronto a schierarsi dalla mia parte in ogni circostanza. Per la nostra amicizia, che è evoluta parecchio rispetto a quando giocavamo con le biglie nella cantina di tua nonna.

A \textbf{Calo}, per la sensibilità che metti sempre in gioco nella nostra amicizia, e per impegnarti costantemente a fare capire quanto ci tieni. Per esserci quando c’è bisogno, e prendermi in giro quando gonfio troppo il petto.

A \textbf{Nicola}, per il carisma e la positività con cui sai affrontare ogni situazione, sorridendo anche quando non è poi così facile farlo.

A \textbf{Nico e Pi}. Non mi è facile descrivere la nostra amicizia, e a pensarci bene non mi è chiaro neanche perché siamo diventati così amici. Grazie perché più di una volta, senza neanche accorgervene, mi avete aiutato a rimettere a fuoco cosa conta davvero quando la bussola non funzionava tanto bene.

Ad \textbf{Anna}, per spingermi costantemente ad andare oltre il mio orticello e per l’empatia con cui affronti gli altri e te stessa. Per essere mia tifosa in ogni situazione e per la grazia con cui non ti siedi mai in prima fila.

A \textbf{Gardo}, per la nostra amicizia che non smette mai di evolvere. Ti ammiro per la tua intelligenza e la persona che sei, ma per me rimarrai sempre quel mediano col codino e i tacchetti di ferro, taglia 47.

A \textbf{Grillo}. Se guardo indietro, i mesi di convivenza sono stati quasi una parentesi onirica della mia vita. Ma che cosa facevamo tutto il giorno? Ho ricordi solo di weekend interi chiusi in casa ad ascoltare la tua playlist, bevendo birra e scrivendo la tesi. Ti ammiro per la tua curiosità infinita e per la lucidità con cui affronti ogni discorso, autentico generatore di idee e riflessioni. Ringrazio anche te Claudio, amico di Federico, che gli stai leggendo questo ringraziamento mentre lanci simulazioni e gli prepari il caffè senza zucchero.

A \textbf{Ila}, per essere la più cara delle mie amiche ed esserci sempre in qualsiasi momento per un caffettino, anche quando le ore di sonno scarseggiano.

A \textbf{Gio Gardo}, per il modo con cui sai trasmettere affetto alle persone a cui vuoi bene.

A \textbf{Pie e Gio Dalpa}, per avermi fatto sempre sentire a casa, mai fuori posto, trattandomi da amico vero, e mai da ragazzino a cui insegnare qualcosa.

Ad \textbf{Edo}, per la tua costante voglia di essere migliore rispetto a ieri e per l’entusiasmo con cui affronti qualsiasi routine.

Ad \textbf{Achille, Lucia e Valentina}, per un’amicizia che continua a crescere di intensità negli anni, nonostante i chilometri che ci separano, e i binari differenti su cui viaggiano le nostre vite.

A \textbf{Leo, Alice e Martina}, per continuare a coltivare la nostra amicizia, anche quando mi dimentico di rispondere al telefono. Ogni volta che ci rivediamo è come se non fossi mai partito da Bologna, e stessimo ancora aspettando il treno degli 05.

A \textbf{Lorenz, Manu e Nicholas}, perché la nostra amicizia è ancora fedele a quel gruppetto di liceali che non si presentava alle interrogazioni di fisica e copiava le versioni di latino. Anche se alcune cose \textit{``siamo troppo vecchi per continuare a farle''}.

Ad \textbf{Alberto}, per il tuo modo totale di vivere l’amicizia, per la costante positività nei rapporti con ogni singola persona e per la tua perenne apertura al nuovo.

A \textbf{Tancredi}, per la tua straordinaria sensibilità e acutezza: incredibile quanto le prime impressioni possano sbagliare. I francesi saranno sempre più forti di te in statistica matematica, ma perlomeno tu sai come si legge $\varphi\iota\lambda\iota\alpha$.

A \textbf{Ross}, per avermi fatto sentire sempre a casa, anche sotto la neve e la perenne pioggia elvetica, soprattutto quando Castel Bolognese sembrava lontano anni luce. E per avermi insegnato a scommettere su se stessi.

A \textbf{Rick e Bond}, per avermi mostrato cosa significa avere vera passione per la materia ed entusiasmo per ogni singolo progetto, senza mai perdere di vista il mondo reale là fuori (e per avermi insegnato a usare il terminale).

A \textbf{Lorenz}, per la tua bontà ed il tuo altruismo, ma altrettanto per l’acume con cui sai mettere in dubbio ogni certezza, costringendomi a non accontentarmi mai delle risposte banali.

A \textbf{Taro} per i mille viaggi in autobus condivisi durante il primo anno, dove ho scoperto che anche dietro al più intelligente degli ingegneri c’è una persona gentile e sensibile.

A \textbf{Ginevra}, per l’attenzione che riservi alle amicizie, e per avermi salvato la carriera insegnandomi che il CV va fatto con UNA SOLA PAGINA.

A \textbf{Justine}, per la tua genuinità e per l’amore che trasmetti per quello che fai… e per i pranzi incastrati (non si sa come) tra mille impegni.

A \textbf{Giuseppe}, per avermi mostrato una nuova prospettiva della vita a Zurigo, e che la fredda città elvetica non è solo ETH Zurigo.

To \textbf{Steven}, for your infinite energy and your continuous hunger for new experiences and for always going out of your way to include me.

To \textbf{Otto}, for showing me how to make the most of life through all the ups and downs, and for reminding me that there’s so much more to it than just work and Switzerland.

To \textbf{Kilian}, for being a constant inspiration to improve myself, for your hunger to build meaningful things, and for the kindness of your person.

To \textbf{Professor Held}, for your continuous trust, mutual respect, and for teaching me that making things complicated is easy, but making simple things that actually work is the real challenge.

\vspace{3em}

Credo di aver finito. (Questa volta la pagina non è bianca).

}




\newpage
% optional, if needed:
\subsection*{Code and data availability}
All analyses in this thesis were carried out in \texttt{R} version 4.5.0 \citep{rcore2025}, building on the \texttt{RBesT} package \citep{weber_applying_2021}. The code that reproduces every figure, table and number reported here is available at \url{https://github.com/SaveFonta/ETH-Statistics-Msc-Thesis}. The dataset for the Crohn's disease case study is the \texttt{crohn} dataset distributed with \texttt{RBesT}.

%%% Local Variables:
%%% mode: latex
%%% TeX-master: "MasterThesisSfS"
%%% End:
